\section{The radial model and the angular obstruction}\label{sec:model}

We first isolate the operation that will change the end. On a round
sphere, the mean curvature and the tangential trace of $K$ are distinct
quantities. Their difference of squares determines a local mass. This
permits us to lower the tensor trace without prescribing that the future
null expansion should vanish.

Fix $m>0$ and a smooth function $v=v(r)$. On any interval where
$1+v^2>2m/r$, define
\begin{equation}\label{eq:radial-model}
 \begin{gathered}
 u^2=1+v^2-2m/r,\qquad
 G=u^{-2}\mathrm dr^2+r^2\sigma,\qquad N=u\partial_r,\\
 k(N,N)=v',\qquad k(N,\cdot)|_{T\Sph}=0,\qquad
 k_T=(v/r)r^2\sigma.
 \end{gathered}
\end{equation}
The positivity condition holds, for example, on $r>2m$.
The leaves have mean curvature $H=2u/r$ and tensor trace $p=2v/r$.
Thus
\[
 \frac r2(1+v^2-u^2)=m,\qquad \theta_+=2(u+v)/r.
\]
Choosing $v=r$ gives the hyperboloidal scale, whereas $v=0$ gives
the spatial Schwarzschild metric. The quantity reduced to zero is
$p=2v/r$, not $\theta_+$.

These data are vacuum for every such choice of $v$. Indeed the
warped-product scalar-curvature calculation gives
\[
 \frac{R_G}{2}=\frac{1-u^2}{r^2}-\frac{2uu'}r,
 \qquad
 \frac{(\tr_Gk)^2-|k|_G^2}{2}
       =\frac{2vv'}r+\frac{v^2}{r^2}.
\]
Their sum is zero by differentiating $u^2=1+v^2-2m/r$.
The radial momentum equals
\[
 \frac{2u}r\left(v'-\frac vr\right)
           -u\partial_r\left(\frac{2v}r\right)=0,
\]
and the angular momentum vanishes by spherical symmetry. This is the
basic flexibility of the construction: changing $v$ alone need not cost
any mass.

For general CKS data the corresponding mass is a function of direction.
If $m$ in \eqref{eq:radial-model} is replaced by $F(r,\omega)$, then
angular derivatives of the lapse appear in both constraints. The energy
contains a leading multiple of $\Delta_\sigma F$ and the angular
momentum a leading multiple of $\mathrm d_{\Sph}F$. We will evolve
the angular function by heat flow and add a trace-free tangential tensor
to cancel these terms. The elementary obstruction to the latter step
explains why balancing belongs at the beginning of the proof.
For a smooth sphere function $f$, write
$\langle f\rangle=\avg{f}$.

\begin{lemma}[The first-harmonic obstruction]\label{lem:obstruction}
If a smooth trace-free symmetric tensor $T$ on the unit sphere satisfies
$\diver_\sigma T=\mathrm d_{\Sph}f$, then
$\int_{\Sph}fn^i\,\mathrm dA_\sigma=0$ for $i=1,2,3$.
Conversely, every smooth $f$ satisfying these three conditions admits
such a tensor.
\end{lemma}
\begin{proof}
The coordinate functions satisfy
$\nabla_\sigma^2n^i=-n^i\sigma$ and $\Delta_\sigma n^i=-2n^i$.
Integration by parts gives
\[
 2\int_{\Sph}fn^i
   =\int_{\Sph}\langle\mathrm df,\mathrm dn^i\rangle
   =-\int_{\Sph}T:\nabla^2n^i=0.
\]
For the converse, write $f=\langle f\rangle+\sum_{\ell\geq2}f_\ell$
in spherical harmonics and let $\lambda_\ell=\ell(\ell+1)$.
The function
\[
 \psi=\sum_{\ell\geq2}\frac{f_\ell}{1-\lambda_\ell/2}
\]
is smooth: repeated integration by parts makes the harmonic
coefficients decay faster than any fixed power of $\lambda_\ell$.
On the unit sphere, commuting derivatives gives
\[
 \diver_\sigma(\nabla^2\psi)^{\TF}
       =\mathrm d_{\Sph}(\tfrac12\Delta_\sigma\psi+\psi)
       =\mathrm d_{\Sph}f.
\]
Thus $T=(\nabla^2\psi)^{\TF}$ works. The constant mode contributes
no gradient; only degree one obstructs the inversion.
\end{proof}

In a balanced chart $f_*=M/4$ has mean $m_B$ and vanishing first
harmonic. We next prove the existence of that chart and bring the
original geometry to round leaves. The small increasing mass correction
used in that collar will continue through the bending. It supplies a
positive DEC margin after the two leading angular terms have canceled.
The round-leaf radial-lapse metric is a zero-shift instance of Bartnik's
quasi-spherical framework \cite{Bartnik,BartnikICM}. Here both the metric
and the second tensor vary, and both constraint inequalities are proved
directly rather than imported from a prescribed-scalar-curvature theorem.
